% Fichier complété par tools/complete_missing_corrections.py.
% Source : DYN/DYN-06-PFD/50_BancBalafre/50_BancBalafre.tex
% Statut : rédigé (4 question(s)).
\begin{corrigebox}[Corrigé rédigé]
\CorrigeQuestion{1}
Le coeur de butée est en translation rectiligne suivant l'axe du banc;
                il ne tourne pas par rapport à 0.
\CorrigeQuestion{2}
Les deux actionneurs symétriques d'un même plan fournissent une
                résultante $2F_V\vec e$ (ou $2F_R\vec e$). Au point $A_4$ :
                $\{T_{V\to CB}\}=\{2F_V\vec e;\vec0\}_{A_4}$ et de même
                $\{T_{R\to CB}\}=\{2F_R\vec e;\vec0\}_{A_8}$, avant transport au point
                commun choisi.
\CorrigeQuestion{3}
On isole $CB$, on transporte les deux torseurs d'action au centre
                d'inertie, puis on applique la résultante dynamique suivant l'axe et le
                moment dynamique autour des axes transverses. La résultante donne
                $2F_V+2F_R-F_{\rm res}=M_{CB}\dot v$; l'équation de moment répartit les
                efforts entre plans avant et arrière selon leurs bras de levier.
\CorrigeQuestion{4}
Le plan dont le bras de levier par rapport au centre d'inertie est le
                plus faible doit fournir l'effort le plus élevé. Le rapport se déduit de
                l'équilibre des moments : $F_V/F_R=d_R/d_V$.
\end{corrigebox}
